\documentclass[a4paper,oneside,12pt]{report}
% Package yang diperlukan
\usepackage{graphicx}      % Untuk memasukkan gambar
\usepackage{titling}       % Untuk kustomisasi halaman judul
\usepackage{geometry}      % Untuk mengatur margin halaman
\usepackage{listings}      % Untuk menampilkan blok kode
\usepackage{xcolor}        % Untuk menggunakan warna
\usepackage{float}         % Untuk kontrol penempatan float (seperti gambar)
\usepackage{hyperref}      % Untuk membuat hyperlink (URL, referensi)
\usepackage{indentfirst}   % Untuk memberi indentasi pada paragraf pertama setiap seksi
\usepackage{url}           % Untuk format URL
\usepackage{fancyhdr}      % Untuk kustomisasi header dan footer
\usepackage{tocloft}       % Untuk kustomisasi Daftar Isi
\usepackage{bookmark}      % Untuk manajemen bookmark PDF yang lebih baik
\usepackage{amsmath}       % Untuk notasi matematika
\usepackage[most]{tcolorbox} % Untuk box output terminal

% Paket untuk tabel agar sesuai format referensi
\usepackage{array}
\usepackage{booktabs}
\usepackage{tabularx}
\usepackage{multirow}
\usepackage{colortbl}

% Macro untuk header tabel
\newcommand{\thead}[1]{\textcolor{white}{\textbf{#1}}}

%======================================================================
% PENGATURAN GEOMETRI HALAMAN
%======================================================================
\geometry{left=3cm, right=2.5cm, top=2cm, bottom=2.5cm}
\sloppy

%======================================================================
% PENGATURAN WARNA UNTUK BLOK KODE
%======================================================================
\definecolor{codegreen}{rgb}{0,0.6,0}
\definecolor{codegray}{rgb}{0.5,0.5,0.5}
\definecolor{codepurple}{rgb}{0.58,0,0.82}
\definecolor{backcolour}{rgb}{0.95,0.95,0.92}
\definecolor{darkblue}{rgb}{0,0,0.6}
\definecolor{outputbg}{RGB}{241,245,249}
\definecolor{outputborder}{RGB}{203,213,225}

%======================================================================
% PENGATURAN BLOK KODE (LISTINGS)
%======================================================================
\lstdefinestyle{commonstyle}{
    backgroundcolor=\color{backcolour},
    breakatwhitespace=false,
    breaklines=true,
    captionpos=b,
    keepspaces=true,
    numbers=left,
    numbersep=5pt,
    showspaces=false,
    showstringspaces=false,
    showtabs=false,
    tabsize=2,
    frame=single,
    basicstyle=\ttfamily\footnotesize,
    numberstyle=\tiny\color{codegray},
    columns=flexible,
    postbreak=\mbox{\textcolor{red}{$\hookrightarrow$}\space},
    breakindent=0pt,
    breakautoindent=false
}

\lstdefinestyle{pythonstyle}{
    style=commonstyle,
    language=Python,
    commentstyle=\color{codegreen},
    keywordstyle=\color{darkblue},
    stringstyle=\color{codepurple},
    emph={import,from,class,def,for,while,if,else,elif,try,except,finally,as,lambda,return},
    emphstyle=\color{darkblue}
}

% ── Box Format Output / Terminal Box (Verbatim via tcblisting) ───────────
\newtcblisting{outputbox}{
    enhanced,
    colback=outputbg,
    colframe=outputborder,
    boxrule=0.6pt,
    arc=1.2mm,
    left=3mm,
    right=3mm,
    top=1.2mm,
    bottom=1.2mm,
    listing only,
    listing options={
        basicstyle=\ttfamily\footnotesize,
        breaklines=true,
        numbers=none,
        showstringspaces=false
    }
}

\newcommand{\screenshotimage}[3]{%
    \begin{figure}[H]
        \centering
        \includegraphics[width=0.85\textwidth]{#1}
        \caption{#2}
        \label{fig:#3}
    \end{figure}
}

%======================================================================
% PENGATURAN LAIN-LAIN
%======================================================================
\renewcommand{\figurename}{Gambar}
\renewcommand{\thefigure}{\arabic{figure}}

\hypersetup{
    colorlinks=true,
    linkcolor=black,
    filecolor=black,
    urlcolor=black,
    citecolor=black,
    pdftitle={Laporan Praktikum Penambangan Data - Pertemuan II},
    pdfpagemode=FullScreen,
    breaklinks=true
}

\pagestyle{fancy}
\fancyhf{}
\fancyfoot[C]{\thepage}
\renewcommand{\headrulewidth}{0pt}

\fancypagestyle{plain}{
    \fancyhf{}
    \fancyfoot[C]{\thepage}
    \renewcommand{\headrulewidth}{0pt}
}

%======================================================================
% PENGATURAN STRUKTUR BAB DAN DAFTAR ISI
%======================================================================
\renewcommand{\contentsname}{Daftar Isi}
\renewcommand{\bibname}{Daftar Pustaka}

\renewcommand{\thechapter}{}
\renewcommand{\chaptername}{}

\renewcommand{\thesection}{\arabic{chapter}.\arabic{section}}
\renewcommand{\thesubsection}{\thesection.\arabic{subsection}}

\setcounter{secnumdepth}{3}
\setcounter{tocdepth}{3}

\renewcommand{\cftchapdotsep}{\cftdotsep}
\renewcommand{\cftchapleader}{\cftdotfill{\cftchapdotsep}}
\setlength{\cftchapindent}{0em}
\setlength{\cftsecindent}{2em}
\setlength{\cftsubsecindent}{4em}
\setlength{\cftchapnumwidth}{0em}
\setlength{\cftsecnumwidth}{2.5em}
\setlength{\cftsubsecnumwidth}{3.2em}
\renewcommand{\cftchappresnum}{}
\renewcommand{\cftchapaftersnum}{}

%======================================================================
% AWAL DOKUMEN
%======================================================================
\begin{document}

\newgeometry{left=3cm, right=2.5cm, top=2.5cm, bottom=3cm}
\begin{titlingpage}
\begin{center}
\vspace*{0.5cm}
{\large \textbf{\MakeUppercase{Laporan Praktikum}}} \\[0.4cm]
{\large \textbf{\MakeUppercase{Praktikum Penambangan Data}}} \\[0.4cm]
{\large \textbf{\MakeUppercase{Pertemuan II}}} \\[0.4cm]
{\large \textbf{\MakeUppercase{Exploratory Data Analysis dan Data Preprocessing}}} \\[1.5cm]

\includegraphics[height=7cm]{lambang ugm.png} \\[1cm]

{\normalsize \textbf{Disusun oleh:}} \\[0.4cm]
\begin{tabular}{@{}l l @{}}
  \textbf{Nama}           & : Januarsyah Akbar \\
  \textbf{NIM}            & : 24/535846/SV/24314 \\
  \textbf{Kelas}          & : B2 \\
  \textbf{Dosen Pengampu} & : Dr. Ganjar Alfian, S.T., M.Eng. \\
                          & \phantom{:} Dr. Imam Fahrurrozi, S.T., M.Cs.
\end{tabular} \\[1.2cm]

{\normalsize \textbf{PROGRAM STUDI D-IV}} \\[0.2cm]
{\normalsize \textbf{TEKNOLOGI REKAYASA PERANGKAT LUNAK}} \\[0.2cm]
{\normalsize \textbf{DEPARTEMEN TEKNIK ELEKTRO DAN INFORMATIKA}} \\[0.2cm]
{\normalsize \textbf{SEKOLAH VOKASI}} \\[0.2cm]
{\normalsize \textbf{UNIVERSITAS GADJAH MADA}} \\[0.2cm]
{\normalsize \textbf{YOGYAKARTA}} \\[0.2cm]
{\normalsize \textbf{2026}} \\[1.2cm]
\end{center}
\end{titlingpage}
\restoregeometry

\thispagestyle{empty}
\clearpage
\stepcounter{page}
\thispagestyle{empty}

% --- Daftar Isi ---
\addcontentsline{toc}{chapter}{Daftar Isi}
\tableofcontents
\clearpage

\pagenumbering{arabic}
\setcounter{page}{3}

%======================================================================
% BAB I: TUJUAN PRAKTIKUM
%======================================================================
\chapter[Tujuan Praktikum]{Tujuan Praktikum}
Tujuan pelaksanaan praktikum pada Pertemuan II adalah sebagai berikut:
\begin{enumerate}
    \item Mahasiswa mampu memahami konsep dasar Artificial Intelligence (AI) dan peran data dalam membangun sistem cerdas.
    \item Mahasiswa mengetahui dasar-dasar Machine Learning (ML) meliputi \textit{supervised learning}, \textit{unsupervised learning}, dan \textit{reinforcement learning}.
    \item Mahasiswa dapat menerapkan langkah awal analisis data melalui \textit{Exploratory Data Analysis} (EDA) menggunakan visualisasi grafik dan eksplorasi statistik.
    \item Mahasiswa mampu mengetahui dan menerapkan teknik-teknik dasar dalam \textit{data preprocessing} meliputi penanganan \textit{missing values}, \textit{duplicate removal}, \textit{categorical encoding}, \textit{dataset splitting}, dan \textit{feature scaling}.
\end{enumerate}
\clearpage

%======================================================================
% BAB II: DASAR TEORI
%======================================================================
\chapter{Dasar Teori}

\section{Artificial Intelligence (AI)}
\textit{Artificial Intelligence} (AI) merupakan cabang ilmu komputer yang berfokus pada pembuatan mesin dan sistem cerdas yang mampu melakukan tugas-tugas yang secara tradisional membutuhkan tingkat kecerdasan manusia. Tugas-tugas ini meliputi penalaran logis, pembelajaran dari pengalaman, pengenalan pola visual maupun bahasa, hingga pengambilan keputusan secara otonom \cite{ref1}. Perkembangan AI modern sangat bertumpu pada ketersediaan data dalam skala besar (\textit{big data}) serta peningkatan kapasitas komputasi perangkat keras seperti GPU dan TPU. Kualitas dari kecerdasan buatan sangat bergantung pada kualitas data yang dipelajari; data yang buruk atau bias akan menghasilkan keputusan sistem yang keliru.

\section{Machine Learning (ML)}
\textit{Machine Learning} adalah subdisiplin utama dari AI yang memungkinkan sistem komputer untuk belajar dan meningkatkan performanya secara otomatis dari pengalaman (data) tanpa diprogram secara eksplisit berbasis aturan kaku (\textit{hardcoded rules}). Secara umum, algoritma \textit{Machine Learning} dikategorikan menjadi tiga paradigma utama \cite{ref2}:
\begin{enumerate}
    \item \textbf{Supervised Learning}: Model dilatih menggunakan dataset terlabel (\textit{labeled data}) di mana terdapat pasangan fitur masukan ($X$) dan target luaran ($y$). Contohnya adalah estimasi harga rumah (regresi) dan deteksi \textit{churn} pelanggan (klasifikasi).
    \item \textbf{Unsupervised Learning}: Model menerima data tanpa label dan bertugas menemukan pola, struktur tersembunyi, atau pengelompokan alami dalam data. Contohnya adalah klusterisasi segmen pelanggan (\textit{K-Means}).
    \item \textbf{Reinforcement Learning}: Agen cerdas belajar mengambil tindakan optimal melalui mekanisme penghargaan (\textit{reward}) dan hukuman (\textit{penalty}) dalam suatu lingkungan dinamis.
\end{enumerate}

\section{Exploratory Data Analysis (EDA)}
\textit{Exploratory Data Analysis} (EDA) adalah pendekatan kritis awal dalam analisis data yang pertama kali dipelopori oleh John W. Tukey. EDA menekankan pentingnya mendengarkan cerita yang disampaikan oleh data sebelum terburu-buru menerapkan model matematika atau algoritma pemodelan yang rumit \cite{ref3,ref5}. Tujuan utama EDA adalah:
\begin{itemize}
    \item Memahami karakteristik dasar dan distribusi dari setiap variabel numerik maupun kategorikal.
    \item Memeriksa tingkat pemusatan data (\textit{mean}, \textit{median}) dan penyebarannya (\textit{standard deviation}, \textit{interquartile range}).
    \item Mengidentifikasi kejanggalan pada data seperti \textit{outliers}, \textit{missing values}, dan rekaman duplikat.
    \item Memetakan korelasi dan hubungan antar variabel menggunakan visualisasi grafik seperti \textit{scatter plot}, \textit{heatmap}, \textit{histogram}, dan \textit{boxplot}.
\end{itemize}

\section{Data Preprocessing}
\textit{Data Preprocessing} adalah serangkaian tahapan sistematis untuk membersihkan, mentransformasikan, dan menyiapkan data mentah agar siap diproses oleh algoritma \textit{Machine Learning}. Data mentah di dunia nyata umumnya bersifat kotor, tidak lengkap, tidak konsisten, serta mengandung \textit{noise}. Tahapan vital dalam \textit{data preprocessing} meliputi \cite{ref4,ref6}:
\begin{enumerate}
    \item \textbf{Data Cleaning}: Penanganan nilai hilang (\textit{imputation}) menggunakan nilai statistik seperti \textit{mean}, \textit{median}, atau modus, serta penghapusan data duplikat (\textit{drop duplicates}).
    \item \textbf{Categorical Encoding}: Mengubah variabel kategorikal teks menjadi representasi angka numerik. Teknik utama meliputi \textit{One-Hot Encoding} (membuat kolom biner baru) dan \textit{Label Encoding} (memetakan kategori ke nilai integer).
    \item \textbf{Dataset Splitting}: Membagi dataset menjadi \textit{training set} (untuk melatih model) dan \textit{testing set} (untuk evaluasi mandiri) guna mencegah fenomena \textit{overfitting}.
    \item \textbf{Feature Scaling}: Menyeleraskan skala variabel numerik agar berada pada rentang yang sebanding. Dua teknik populer adalah \textit{Standardization} (z-score scaling dengan $\mu=0, \sigma=1$) dan \textit{Normalization} (Min-Max scaling pada rentang $[0, 1]$).
\end{enumerate}
\clearpage

%======================================================================
% BAB III: HASIL DAN PEMBAHASAN
%======================================================================
\chapter{Hasil dan Pembahasan}

\section{Langkah Percobaan}

\subsection{Langkah Percobaan Exploratory Data Analysis (EDA)}
Pada bagian ini, dilakukan serangkaian eksplorasi awal terhadap dataset \textit{BankChurners} dan data sensor IoT untuk mengenali bentuk sebaran datanya:

\begin{enumerate}
    \item \textbf{Impor Pustaka Utama} \\
    Pertama-tama, disiapkan beberapa pustaka Python standar yang umum digunakan untuk pengolahan data dan visualisasi, yaitu \texttt{pandas}, \texttt{matplotlib.pyplot}, dan \texttt{numpy}.
\begin{lstlisting}[style=pythonstyle, caption=Impor Pustaka EDA]
import pandas as pd
import matplotlib.pyplot as plt
import numpy as np
\end{lstlisting}

    \item \textbf{Load Dataset Bank Churners} \\
    Dataset dibaca langsung dari tautan repositori daring menggunakan fungsi \texttt{pd.read\_csv()}.
\begin{lstlisting}[style=pythonstyle, caption=Load Dataset Bank Churners]
df_churning = pd.read_csv('https://raw.githubusercontent.com/ganjar87/data_science_practice/main/BankChurners.csv')
print('Shape:', df_churning.shape)
\end{lstlisting}
    \screenshotimage{gambar/out_churn_shape.png}{Output Dimensi Baris dan Kolom Dataset Bank Churners}{out_churn_shape}
    Dari hasil pembacaan di atas, terlihat bahwa dataset memiliki ukuran sebanyak 10.127 baris dan 23 kolom fitur.

    \item \textbf{Melihat 5 Baris Pertama dan Terakhir Data} \\
    Untuk memperoleh gambaran awal struktur data, dilakukan inspeksi 5 baris teratas menggunakan \texttt{head()} dan 5 baris terbawah menggunakan \texttt{tail()}.
\begin{lstlisting}[style=pythonstyle, caption=Inspeksi Baris Pertama dan Terakhir]
df_churning.head()
df_churning.tail()
\end{lstlisting}

    \item \textbf{Pemisahan Sub-Dataset Berdasarkan Kelas Target} \\
    Agar analisis perbandingan lebih mudah dilakukan, dataset dipisahkan menjadi dua subset berdasarkan status nasabah pada kolom \texttt{Attrition\_Flag}, yaitu kelompok nasabah yang masih aktif (\textit{existing}) dan nasabah yang sudah berhenti (\textit{attrited}).
\begin{lstlisting}[style=pythonstyle, caption=Pemisahan Data Existing vs Attrited]
existing_data = df_churning[(df_churning.Attrition_Flag == 'Existing Customer')]
attrited_data = df_churning[(df_churning.Attrition_Flag == 'Attrited Customer')]
\end{lstlisting}

    \item \textbf{Line Chart (Data Dummy \& IoT Humidity Time Series)} \\
    Grafik garis ini digunakan untuk mengamati pola fluktuasi kelembaban udara dari sensor IoT seiring berjalannya waktu.
\begin{lstlisting}[style=pythonstyle, caption=Visualisasi Line Chart IoT]
df_iot = pd.read_csv('https://raw.githubusercontent.com/ganjar87/data_science_practice/main/datatraining.txt')
df_iot['date'] = pd.to_datetime(df_iot['date'])
plt.figure(figsize=(15, 3))
plt.plot(df_iot['date'], df_iot['Humidity'])
plt.show()
\end{lstlisting}
    \screenshotimage{gambar/fig_iot_humidity.png}{Grafik Line Chart Tren Kelembaban Sensor IoT}{iot_humidity}
    Dapat diamati pada grafik garis di atas, nilai kelembaban (\textit{humidity}) dari sensor IoT tampak berosilasi secara periodik mengikuti pola waktu tertentu (\textit{time series}). Pola siklikal harian ini wajar terjadi seiring pergantian siang dan malam di sekitar lingkungan sensor.

    \item \textbf{Pie Chart (Proporsi Marital Status)} \\
    Diagram lingkaran ini diterapkan untuk melihat persentase perbandingan status pernikahan nasabah kartu kredit.
\begin{lstlisting}[style=pythonstyle, caption=Visualisasi Pie Chart Marital Status]
data = df_churning['Marital_Status'].value_counts()
plt.pie(data, labels=data.index, autopct='%.2f%%')
plt.title('Marital Status')
plt.show()
\end{lstlisting}
    \screenshotimage{gambar/fig_pie_marital.png}{Pie Chart Distribusi Marital Status Nasabah}{pie_marital}
    Berdasarkan diagram lingkaran pada Gambar \ref{fig:pie_marital}, dapat diketahui bahwa sebagian besar nasabah telah berstatus menikah (\textit{Married}) yaitu sebesar 46,28\%, disusul nasabah lajang (\textit{Single}) sebesar 38,94\%. Sisanya merupakan nasabah dengan status belum tercatat (\textit{Unknown}) sekitar 7,40\% dan nasabah bercerai (\textit{Divorced}) sebesar 7,39\%.

    \item \textbf{Grouped Bar Plot \& Stacked Bar Plot} \\
    Visualisasi diagram batang (berdampingan dan bertumpuk) digunakan untuk membandingkan komposisi status pernikahan antara nasabah aktif dan nasabah yang mengalami \textit{churn}.
\begin{lstlisting}[style=pythonstyle, caption=Grouped dan Stacked Bar Plot]
# Stacked Bar Plot
data1 = existing_data['Marital_Status'].value_counts()
data2 = attrited_data['Marital_Status'].value_counts()
df_new = pd.concat([data1, data2], keys=['existing', 'attrited'], axis=1)
df_new.plot(kind='bar', stacked=True)
plt.title('Marital Status')
plt.ylabel('Count')
plt.show()
\end{lstlisting}
    \screenshotimage{gambar/fig_bar_marital_grouped.png}{Grouped Bar Plot Status Pernikahan per Kelompok Nasabah}{bar_grouped}
    \screenshotimage{gambar/fig_bar_marital_stacked.png}{Stacked Bar Plot Status Pernikahan Nasabah}{bar_stacked}
    Berdasarkan tampilan \textit{grouped} dan \textit{stacked bar plot} di atas, terlihat bahwa proporsi nasabah yang mengalami \textit{churn} tersebar merata di setiap kategori status pernikahan. Kategori \textit{Married} dan \textit{Single} tetap menjadi kelompok dengan jumlah nasabah paling dominan pada kedua kelompok nasabah.

    \item \textbf{Histogram Distribusi Usia (Customer Age)} \\
    Histogram ini dirancang untuk melihat pola sebaran rentang usia nasabah aktif dibandingkan dengan nasabah yang \textit{churn}.
\begin{lstlisting}[style=pythonstyle, caption=Histogram Usia Pelanggan]
plt.hist(existing_data['Customer_Age'], label='existing')
plt.hist(attrited_data['Customer_Age'], label='attrited')
plt.title('Customer Age')
plt.legend()
plt.show()
\end{lstlisting}
    \screenshotimage{gambar/fig_hist_age.png}{Histogram Distribusi Usia Nasabah Aktif vs Churn}{hist_age}
    Kurva histogram pada Gambar \ref{fig:hist_age} memperlihatkan bahwa distribusi usia untuk kedua kelompok nasabah membentuk pola lonceng yang simetris (distribusi normal). Kepadatan nasabah tertinggi berada pada rentang usia produktif antara 40 hingga 55 tahun.

    \item \textbf{Scatter Plot (Customer Age vs Months on Book)} \\
    Diagram pencar (\textit{scatter plot}) ini digunakan untuk mengidentifikasi korelasi antara usia nasabah dengan lama masa berlangganan di bank.
\begin{lstlisting}[style=pythonstyle, caption=Scatter Plot Usia vs Masa Langganan]
plt.scatter(existing_data['Customer_Age'], existing_data['Months_on_book'], label='existing')
plt.scatter(attrited_data['Customer_Age'], attrited_data['Months_on_book'], label='attrited')
plt.title('Customer age vs months on book')
plt.show()
\end{lstlisting}
    \screenshotimage{gambar/fig_scatter_age_months.png}{Scatter Plot Usia Nasabah vs Months on Book}{scatter_age}
    Sebaran titik pada Gambar \ref{fig:scatter_age} mengindikasikan kecenderungan hubungan linear positif: semakin bertambah usia nasabah, durasi masa keanggotaan di bank (\textit{Months on Book}) juga cenderung semakin lama.

    \item \textbf{Boxplot \& Violin Plot} \\
    Kedua diagram ini membantu dalam mengevaluasi ringkasan 5 angka statistik (median, kuartil bawah, kuartil atas, nilai minimum, dan maksimum), keberadaan pencilan (\textit{outlier}), serta estimasi kepadatan data (\textit{KDE}).
\begin{lstlisting}[style=pythonstyle, caption=Boxplot dan Violin Plot]
# Violin plot
dict_violin = {'existing': existing_data['Customer_Age'].values, 'attrited': attrited_data['Customer_Age'].values}
plt.violinplot(list(dict_violin.values()))
plt.xticks([1, 2], labels=list(dict_violin.keys()))
plt.show()
\end{lstlisting}
    \screenshotimage{gambar/fig_boxplot_grouped.png}{Boxplot Perbandingan Usia Nasabah}{boxplot_grouped}
    \screenshotimage{gambar/fig_violin_age.png}{Violin Plot Kepadatan Usia Nasabah}{violin_age}
    Melalui \textit{boxplot} dan \textit{violin plot} di atas, dapat diamati bahwa nilai median usia nasabah aktif dan nasabah \textit{churn} relatif identik di kisaran angka 46 tahun. Rentang penyebaran antarkuartil juga tampak simetris tanpa adanya pencilan (\textit{outlier}) yang ekstrem.

    \item \textbf{Subplot 2x2 \& Anotasi Statistik} \\
    Pada tahap ini, dilakukan penggabungan beberapa grafik ke dalam satu bidang kanvas (kisi 2x2) serta penambahan teks anotasi langsung pada grafik untuk mencantumkan nilai rata-rata (\textit{mean}) dan nilai tengah (\textit{median}).
\begin{lstlisting}[style=pythonstyle, caption=Anotasi Ukuran Pemusatan Data]
data = df_churning['Customer_Age']
plt.hist(data)
plt.text(60, 1600, "Median Age: " + str(round(data.median(), 2)), style='italic')
plt.text(60, 1800, "Mean Age: " + str(round(data.mean(), 2)), style='italic')
plt.show()
\end{lstlisting}
    \screenshotimage{gambar/fig_subplot_2x2.png}{Subplot Kisi 2x2 Boxplot dan Histogram Usia}{subplot_2x2}
    \screenshotimage{gambar/fig_annotation.png}{Histogram Usia Dilengkapi Anotasi Statistics}{annotation_plot}
    Pada Gambar \ref{fig:subplot_2x2}, tata letak subplot 2x2 mempermudah perbandingan visual secara terpadu. Sementara itu pada Gambar \ref{fig:annotation_plot}, anotasi statistik memperlihatkan nilai \textit{mean} sebesar 46,33 tahun dan \textit{median} 46,00 tahun. Kedekatan kedua nilai ini mengonfirmasi bahwa distribusi usia nasabah berkarakteristik simetris normal.
\end{enumerate}

\subsection{Langkah Percobaan Data Preprocessing (Titanic Dataset)}
Berikut adalah langkah-langkah penyiapan dan pembersihan data mentah yang diterapkan pada dataset \textit{Titanic}:

\begin{enumerate}
    \item \textbf{Load Dataset \& Cek Dimensi} \\
    Pertama, dataset pelatihan \textit{Titanic} dimuat lalu diperiksa ukuran dimensinya menggunakan fungsi \texttt{shape}.
\begin{lstlisting}[style=pythonstyle, caption=Load Dataset Titanic]
df = pd.read_csv('https://raw.githubusercontent.com/ganjar87/data_science_practice/main/train.csv')
print('Shape:', df.shape)
\end{lstlisting}
    \screenshotimage{gambar/out_titanic_shape.png}{Output Dimensi Baris dan Kolom Dataset Titanic}{out_titanic_shape}
    Dataset Titanic ini tercatat memiliki 891 baris data penumpang dan 12 kolom atribut.

    \item \textbf{Inspeksi Missing Values \& Duplikasi} \\
    Pemeriksaan dilakukan untuk mendeteksi keberadaan data kosong (\textit{missing values}) maupun baris data yang terduplikasi.
\begin{lstlisting}[style=pythonstyle, caption=Pemeriksaan Nilai Kosong dan Duplikat]
print(df.isnull().sum())
print('Duplikat:', df.duplicated().sum())
df.drop_duplicates(inplace=True)
\end{lstlisting}
    \screenshotimage{gambar/out_titanic_null_dup.png}{Output Pemeriksaan Missing Values dan Duplikasi Titanic}{out_titanic_null_dup}
    Berdasarkan output di atas, tidak ditemukan rekaman data duplikat (0 duplikat). Namun demikian, terdapat nilai kosong pada kolom \texttt{Age} (177 data), kolom \texttt{Cabin} (687 data), dan kolom \texttt{Embarked} (2 data) yang memerlukan proses imputasi.

    \item \textbf{Imputasi Nilai Kosong (Imputation)} \\
    Untuk menangani data yang hilang, kolom numerik \texttt{Age} diisi menggunakan nilai tengah (\textit{median}), sedangkan kolom kategorikal seperti \texttt{Cabin} dan \texttt{Embarked} diisi menggunakan nilai dengan frekuensi kemunculan tertinggi (modus).
\begin{lstlisting}[style=pythonstyle, caption=Imputasi Fitur Titanic]
df['Age'].fillna(df['Age'].median(), inplace=True)
df['Cabin'].fillna(df['Cabin'].value_counts().index[0], inplace=True)
df['Embarked'].fillna(df['Embarked'].value_counts().index[0], inplace=True)
\end{lstlisting}

    \item \textbf{Encoding Variabel Kategorikal} \\
    Agar dapat diproses oleh algoritma machine learning, atribut teks kategorikal seperti \texttt{Sex} dan \texttt{Embarked} ditransformasikan menjadi representasi numerik menggunakan \textit{One-Hot Encoding} (\texttt{pd.get\_dummies}) serta \textit{Label Encoding} (\texttt{LabelEncoder}).
\begin{lstlisting}[style=pythonstyle, caption=One-Hot Encoding dan Label Encoding]
df_X = df.drop(['PassengerId', 'Name', 'Survived', 'Cabin', 'Ticket'], axis=1)
df_y = df['Survived']

# One-Hot Encoding
df_onehot = pd.get_dummies(df_X, columns=['Sex', 'Embarked'], drop_first=True)

# Label Encoding
cat_en = LabelEncoder()
for col in ['Sex', 'Embarked']:
    df_X[col] = cat_en.fit_transform(df_X[col])
\end{lstlisting}

    \item \textbf{Dataset Splitting (Train-Test Split)} \\
    Supaya evaluasi model berjalan secara objektif dan terbebas dari kebocoran data, dataset dipisahkan dengan proporsi 70\% untuk data latih (\textit{training set}) dan 30\% untuk data uji (\textit{testing set}).
\begin{lstlisting}[style=pythonstyle, caption=Pembagian Dataset 70:30]
X_train, X_test, y_train, y_test = train_test_split(df_X, df_y, test_size=0.3, random_state=42)
print('X_train shape:', X_train.shape, '| X_test shape:', X_test.shape)
\end{lstlisting}
    \screenshotimage{gambar/out_titanic_split.png}{Output Pembagian Dimensi Training Set dan Testing Set}{out_titanic_split}
    Pembagian ini menghasilkan 623 baris sampel untuk pelatihan dan 268 baris sampel mandiri untuk pengujian model.

    \item \textbf{Feature Scaling (StandardScaler \& MinMaxScaler)} \\
    Terakhir, dilakukan penyelarasan skala fitur numerik agar variabel berskala besar tidak mendominasi proses pembelajaran model. Proses penyesuaian (\textit{fit}) dilakukan **khusus pada data latih saja**, lalu hasilnya ditransformasikan ke data latih dan data uji.
\begin{lstlisting}[style=pythonstyle, caption=Standardization dan Normalization]
# StandardScaler (Z-Score)
scaler_std = StandardScaler()
X_train_std = scaler_std.fit_transform(X_train)
X_test_std = scaler_std.transform(X_test)

# MinMaxScaler ([0, 1])
scaler_mm = MinMaxScaler()
X_train_mm = scaler_mm.fit_transform(X_train)
X_test_mm = scaler_mm.transform(X_test)
\end{lstlisting}
\end{enumerate}

\clearpage

%======================================================================
% TUGAS DAN ANALISIS
%======================================================================
\section{Tugas dan Analisis}

\subsection{Bagian A: Exploratory Data Analysis (Bank Churners Dataset)}
Pada tugas bagian A ini, analisis difokuskan untuk mengeksplorasi dataset \textit{Credit Card Customers} (\texttt{BankChurners.csv}) guna memahami karakteristik nasabah kartu kredit secara mendalam.

\begin{enumerate}
    \item \textbf{Membaca Dataset Menggunakan Dataframe} \\
    Pertama, dataset dimuat ke dalam dataframe Pandas untuk dianalisis dimensi bentuknya.
\begin{lstlisting}[style=pythonstyle, caption=Load Dataset Tugas EDA]
df_tugas_eda = pd.read_csv('https://raw.githubusercontent.com/ganjar87/data_science_practice/main/BankChurners.csv')
print('Dimensi Dataset:', df_tugas_eda.shape)
\end{lstlisting}
    \screenshotimage{gambar/out_tugas_churn_shape.png}{Output Dimensi Dataset Bank Churners}{out_tugas_churn_shape}

    \item \textbf{Tujuan Penggunaan Dataset} \\
    Secara umum, dataset \textit{BankChurners} ini digunakan untuk menganalisis dan memprediksi kemungkinan nasabah dalam menghentikan pemakaian fasilitas kartu kredit (\textit{customer churn}). Melalui dataset ini, pihak manajemen bank dapat memetakan perbedaan pola antara nasabah yang tetap aktif bertransaksi (\texttt{Existing Customer}) dengan nasabah yang memilih menutup akun (\texttt{Attrited Customer}). Karakteristik ini menjadi dasar penting dalam pengembangan model machine learning klasifikasi terawasi (\textit{supervised learning}) guna mendeteksi nasabah berisiko tinggi secara dini, sehingga program retensi dapat dirancang lebih cepat dan tepat sasaran.

    \item \textbf{Definisi Atribut Input, Output/Label, dan Kolom yang Dibuang} \\
    Dalam pemodelan data ini, target utama yang menjadi variabel luaran (\textit{target label}) adalah kolom \texttt{Attrition\_Flag}. Beberapa kolom perlu dieliminasi dari daftar fitur masukan. Kolom \texttt{CLIENTNUM} dibuang karena murni berfungsi sebagai nomor identitas pencatatan tanpa korelasi terhadap pola churn. Dua kolom bawaan pengklasifikasi Naive Bayes di bagian akhir tabel (\texttt{Naive\_Bayes\_Classifier\_...}) juga wajib dihapus untuk mencegah kebocoran informasi evaluasi (\textit{data leakage}). Sementara itu, 19 kolom sisanya ditetapkan sebagai fitur masukan (\textit{features}) yang mencakup variabel demografi, status kepemilikan akun, riwayat limit dan saldo berputar, frekuensi kontak, hingga rasio volume transaksi.

    \item \textbf{Penjelasan Singkat Setiap Atribut Dataset} \\
    Rincian fungsi dan tipe data dari 21 kolom dasar pada dataset \textit{BankChurners} disajikan pada Tabel~\ref{tab:bank_churners_attributes}.

\begin{table}[H]
\centering
\caption{Deskripsi dan Karakteristik Atribut Dataset BankChurners}
\label{tab:bank_churners_attributes}
\footnotesize
\rowcolors{2}{gray!10}{white}
\begin{tabularx}{\textwidth}{l l X}
\toprule
\rowcolor{darkblue}
\thead{Nama Atribut} & \thead{Tipe Data} & \thead{Penjelasan / Fungsi Atribut} \\
\midrule
\texttt{CLIENTNUM} & Integer & Nomor ID unik pemegang kartu kredit (di-drop). \\
\texttt{Attrition\_Flag} & Object & Target label: \textit{Existing Customer} vs \textit{Attrited Customer}. \\
\texttt{Customer\_Age} & Integer & Usia nasabah dalam satuan tahun. \\
\texttt{Gender} & Object & Jenis kelamin nasabah (\textit{M} = Pria, \textit{F} = Wanita). \\
\texttt{Dependent\_count} & Integer & Jumlah anggota keluarga yang menjadi tanggungan. \\
\texttt{Education\_Level} & Object & Tingkat pendidikan terakhir nasabah. \\
\texttt{Marital\_Status} & Object & Status pernikahan (\textit{Married}, \textit{Single}, \textit{Divorced}, \textit{Unknown}). \\
\texttt{Income\_Category} & Object & Kelompok rentang penghasilan tahunan nasabah. \\
\texttt{Card\_Category} & Object & Tipe tingkatan kartu kredit (\textit{Blue}, \textit{Silver}, \textit{Gold}, \textit{Platinum}). \\
\texttt{Months\_on\_book} & Integer & Lama waktu nasabah menjadi nasabah bank (bulan). \\
\texttt{Total\_Relationship\_Count} & Integer & Total jumlah produk perbankan yang dimiliki nasabah. \\
\texttt{Months\_Inactive\_12\_mon} & Integer & Jumlah bulan tidak aktif transaksi dalam 12 bulan terakhir. \\
\texttt{Contacts\_Count\_12\_mon} & Integer & Frekuensi interaksi/kontak nasabah ke bank dalam 12 bulan. \\
\texttt{Credit\_Limit} & Float & Batas maksimum limit kredit yang disetujui bank. \\
\texttt{Total\_Revolving\_Bal} & Integer & Total saldo berputar (\textit{revolving balance}) pada kartu kredit. \\
\texttt{Avg\_Open\_To\_Buy} & Float & Rata-rata sisa limit kredit yang masih dapat digunakan. \\
\texttt{Total\_Amt\_Chng\_Q4\_Q1} & Float & Rasio perubahan nominal transaksi Q4 dibanding Q1. \\
\texttt{Total\_Trans\_Amt} & Integer & Total nominal rupiah/dolar transaksi dalam 12 bulan. \\
\texttt{Total\_Trans\_Ct} & Integer & Total frekuensi jumlah kali transaksi dalam 12 bulan. \\
\texttt{Total\_Ct\_Chng\_Q4\_Q1} & Float & Rasio perubahan frekuensi transaksi Q4 dibanding Q1. \\
\texttt{Avg\_Utilization\_Ratio} & Float & Rasio pemakaian limit kredit (\textit{saldo} / \textit{limit}). \\
\bottomrule
\end{tabularx}
\end{table}
\end{enumerate}

\subsection{Bagian B: Data Preprocessing (Bengaluru House Price Dataset)}
Pada tugas bagian B ini, tahapan pembersihan dan penyiapan data lengkap diterapkan pada dataset harga perumahan di kota Bengaluru (\textit{Bengaluru House Price}).

\begin{enumerate}
    \item \textbf{Membaca Dataset Menggunakan Dataframe} \\
    Dataset mentah properti perumahan Bengaluru dibaca langsung dari tautan repositori daring.
\begin{lstlisting}[style=pythonstyle, caption=Load Dataset Bengaluru House Price]
URL_BENGALURU = 'https://raw.githubusercontent.com/dphi-official/Datasets/master/Bengaluru_House_Data.csv'
df_house = pd.read_csv(URL_BENGALURU)
print('Dimensi:', df_house.shape)
df_house.head()
\end{lstlisting}
    \screenshotimage{gambar/out_bengaluru_shape.png}{Output Dimensi Baris dan Kolom Dataset Bengaluru}{out_bengaluru_shape}
    Dataset properti ini memiliki catatan sebanyak 13.320 baris rumah dengan 9 kolom atribut.

    \item \textbf{Visualisasi Histogram Numerik \& Analisis Pola Distribusi} \\
    Sebelum pembuatan visualisasi, kolom teks \texttt{total\_sqft} yang berformat rentang (misal \texttt{2100 - 2850}) dikonversi menjadi nilai numerik rata-rata (\texttt{total\_sqft\_num}), serta dilakukan ekstraksi jumlah kamar tidur (\texttt{bhk}) dari kolom \texttt{size}.
\begin{lstlisting}[style=pythonstyle, caption=Parsing Kolom Numerik dan Histogram]
# Visualisasi histogram 5 variabel numerik
num_cols = ['total_sqft_num', 'bath', 'balcony', 'bhk', 'price']
fig, axes = plt.subplots(2, 3, figsize=(10, 6))
# ... [Plotting Histogram] ...
\end{lstlisting}
    \screenshotimage{gambar/fig_bengaluru_hist.png}{Histogram Distribusi Variabel Numerik Dataset Bengaluru}{bengaluru_hist}

    \textit{Penjelasan Pola Distribusi:}
    \begin{itemize}
        \item \textbf{\texttt{price} (Harga Rumah)}: Terlihat pola kemiringan kanan yang tajam (\textit{right-skewed}). Mayoritas harga rumah terkonsentrasi pada segmen menengah ke bawah di bawah 200 Lakh, dengan sedikit sekali properti bernilai sangat tinggi.
        \item \textbf{\texttt{total\_sqft\_num} (Luas Bangunan)}: Menunjukkan distribusi condong ke kanan (\textit{right-skewed}). Sebagian besar bangunan berukuran di bawah 2.000 kaki persegi, disertai keberadaan beberapa pencilan bertanah luas.
        \item \textbf{\texttt{bath} \& \texttt{bhk} (Kamar Mandi \& Kamar Tidur)}: Memiliki sebaran nilai diskrit berskala kecil yang terpusat pada 2 hingga 3 kamar tidur dan kamar mandi per unit rumah.
        \item \textbf{\texttt{balcony} (Balkon)}: Jumlah balkon berada pada rentang 0 hingga 3 unit, dengan frekuensi tertinggi pada unit berfasilitas 1 dan 2 balkon.
    \end{itemize}

    \item \textbf{Matriks Korelasi Antar Variabel \& Analisis Hubungan} \\
    Kalkulasi koefisien korelasi Pearson antar fitur numerik divisualisasikan dalam bentuk \textit{heatmap} untuk mempermudah identifikasi asosiasi antar variabel.
\begin{lstlisting}[style=pythonstyle, caption=Korelasi Heatmap Dataset Bengaluru]
corr_house = df_house[num_cols].corr(numeric_only=True)
# ... [Plotting Heatmap] ...
\end{lstlisting}
    \screenshotimage{gambar/fig_bengaluru_corr.png}{Heatmap Korelasi Pearson Properti Bengaluru}{bengaluru_corr}

    \textit{Analisis Hubungan Antar Variabel:}
    \begin{itemize}
        \item \textbf{\texttt{bath} vs \texttt{bhk} (Korelasi Kuat $\approx 0.90$)}: Hubungan linear positif sangat kuat. Penambahan jumlah kamar tidur secara konsisten diimbangi oleh penambahan fasilitas kamar mandi.
        \item \textbf{\texttt{total\_sqft\_num} vs \texttt{price} (Korelasi Positif $\approx 0.57$)}: Luas total area bangunan merupakan faktor utama yang mendorong kenaikan harga properti di Bengaluru.
        \item \textbf{\texttt{bhk}/\texttt{bath} vs \texttt{price} (Korelasi Sedang $\approx 0.40 - 0.45$)}: Jumlah kamar tidur dan kamar mandi berbanding lurus dengan peningkatan nilai jual rumah.
        \item \textbf{\texttt{balcony} vs \texttt{price} (Korelasi Lemah $\approx 0.12$)}: Jumlah ketersediaan balkon memiliki dampak yang relatif kecil terhadap penentuan harga properti.
    \end{itemize}

    \item \textbf{Exploratory Data Analysis: Statistik Deskriptif, Missing Values, dan Duplikasi} \\
    Pemeriksaan kualitas data dilakukan melalui kalkulasi nilai kosong menggunakan \texttt{isnull().sum()} dan duplikasi baris dengan \texttt{duplicated().sum()}.
\begin{lstlisting}[style=pythonstyle, caption=Pemeriksaan Ringkasan Statistik dan Kualitas Data]
# Missing values check
print(df_house.isnull().sum())
print('Jumlah Duplikat:', df_house.duplicated().sum())
\end{lstlisting}
    \screenshotimage{gambar/out_bengaluru_null_dup.png}{Output Pemeriksaan Missing Values dan Duplikasi Bengaluru}{out_bengaluru_null_dup}

    \textit{Penjelasan Tindakan Penanganan Data:}
    \begin{enumerate}
        \item \textbf{Data Duplikat (568 baris)}: Baris duplikat dihapus menggunakan \texttt{drop\_duplicates()} guna menghindari pembobotan sampel yang berlebihan saat melatih model.
        \item \textbf{Fitur \texttt{society} (5.502 data kosong $\approx 41.3\%$)}: Kolom ini di-drop sepenuhnya karena porsi kekosongan melebihi 40\%, sehingga imputasi buatan justru berisiko merusak integritas data.
        \item \textbf{Fitur \texttt{balcony} (605 data kosong)}: Karena bertipe numerik diskrit dengan sebaran miring, nilai kosong diisi menggunakan nilai tengah (\textit{median}).
        \item \textbf{Fitur \texttt{location}, \texttt{size}, \texttt{bath} (missing $< 1\%$)}: Karena persentase nilai kosong sangat rendah terhadap total dataset, baris yang tidak lengkap pada kolom-kolom ini dieliminasi dengan \texttt{dropna()}.
    \end{enumerate}

    \item \textbf{Preprocessing Lanjutan: Encoding, Split 80:20, dan Feature Scaling} \\
    Setelah tahapan pembersihan selesai, pipeline penyiapan fitur diimplementasikan secara sistematis:
\begin{lstlisting}[style=pythonstyle, caption=Pipeline Preprocessing Lengkap Dataset Bengaluru]
# 1. Cleaning & Imputation
df_pp = df_house.drop_duplicates().drop(columns=['society'])
df_pp = df_pp.dropna(subset=['location', 'size', 'bath', 'bhk', 'total_sqft_num'])
df_pp['balcony'] = df_pp['balcony'].fillna(df_pp['balcony'].median())

# 2. Categorical Encoding
# One-Hot Encoding pada area_type (kardinalitas rendah)
X_ohe = pd.get_dummies(df_pp.drop(columns=['availability', 'size', 'total_sqft', 'price']), 
                       columns=['area_type'], drop_first=True)
# Label Encoding pada location (kardinalitas tinggi > 1000 lokasi)
le_loc = LabelEncoder()
X_ohe['location'] = le_loc.fit_transform(X_ohe['location'])

# 3. Train-Test Split 80:20
y_house = df_pp['price']
X_train_h, X_test_h, y_train_h, y_test_h = train_test_split(
    X_ohe, y_house, test_size=0.2, random_state=42
)

# 4. Feature Scaling (StandardScaler)
num_scale_cols = ['total_sqft_num', 'bath', 'balcony', 'bhk']
scaler_house = StandardScaler()
X_train_scaled = X_train_h.copy()
X_test_scaled = X_test_h.copy()
X_train_scaled[num_scale_cols] = scaler_house.fit_transform(X_train_h[num_scale_cols])
X_test_scaled[num_scale_cols] = scaler_house.transform(X_test_h[num_scale_cols])
\end{lstlisting}

    \textit{Alasan Penggunaan Masing-Masing Teknik:}
    \begin{itemize}
        \item \textbf{One-Hot Encoding (\texttt{area\_type})}: Digunakan untuk tipe area karena hanya memiliki 4 kategori unik, sehingga mencegah timbulnya asumsi urutan/peringkat buatan. Parameter \texttt{drop\_first=True} dipasang untuk menghindari perangkap multikolinearitas (\textit{dummy variable trap}).
        \item \textbf{Label Encoding (\texttt{location})}: Diterapkan pada nama lokasi karena terdapat lebih dari 1.000 lokasi unik. Penggunaan One-Hot pada fitur ini akan memicu ledakan dimensi data (\textit{curse of dimensionality}).
        \item \textbf{Train-Test Split (80:20)}: Membagi 80\% data untuk pelatihan model dan 20\% untuk pengujian objektif, guna memastikan model mampu mengenali pola data baru secara akurat.
        \item \textbf{StandardScaler}: Penyelarasan skala numerik dilakukan agar berpusat pada rata-rata 0 dan variansi 1, yang terbukti lebih stabil dan tahan terhadap keberadaan nilai pencilan (\textit{outlier}) dibanding normalisasi min-max.
    \end{itemize}
\end{enumerate}

\clearpage

%======================================================================
% BAB IV: KESIMPULAN
%======================================================================
\chapter{Kesimpulan}
Berdasarkan seluruh rangkaian praktikum yang telah dilaksanakan pada Pertemuan II, dapat disimpulkan bahwa tahapan \textit{Exploratory Data Analysis} (EDA) dan \textit{Data Preprocessing} memegang peranan yang sangat krusial dalam menentukan keberhasilan pembuatan model \textit{Machine Learning}. Melalui EDA, karakteristik alami data dapat dikenali, pola distribusi dapat dipahami, serta korelasi antar variabel dapat dipetakan secara visual. Pemilihan grafik yang tepat—mulai dari \textit{line chart} untuk tren waktu, diagram batang dan lingkaran untuk proporsi kategori, hingga kombinasi histogram, \textit{scatter plot}, \textit{boxplot}, dan \textit{violin plot} untuk data numerik—sangat membantu dalam mengidentifikasi anomali seperti data pencilan (\textit{outliers}), nilai kosong, dan duplikasi sejak awal.

Selain memahami kondisi data, tahapan \textit{Data Preprocessing} menjadi pondasi utama untuk menjamin data mentah siap diproses oleh algoritma pemodelan secara optimal. Pembersihan data harus dilakukan dengan cermat dan disesuaikan dengan konteks fitur, seperti mengisi nilai kosong dengan median untuk data numerik miring atau modus untuk data kategorikal, serta membuang kolom yang memiliki nilai kosong terlalu dominan. Begitu pula dalam pengkodean variabel kategorikal, pemisahan metode antara \textit{One-Hot Encoding} untuk kategori sedikit dan \textit{Label Encoding} untuk kategori berkardinalitas tinggi diperlukan guna mencegah bias urutan maupun lonjakan dimensi. Terakhir, proses pembagian dataset (\textit{Train-Test Split}) sebelum penerapan penyeimbangan skala (\textit{Feature Scaling}) terbukti sangat penting untuk mencegah terjadinya kebocoran informasi (\textit{data leakage}), sehingga model yang dibangun memiliki kemampuan generalisasi yang optimal saat dievaluasi pada data baru.

\clearpage

%======================================================================
% DAFTAR PUSTAKA
%======================================================================
\begin{thebibliography}{99}
\addcontentsline{toc}{chapter}{Daftar Pustaka}

\bibitem{ref1}
Russell, S., \& Norvig, P. (2020). \textit{Artificial Intelligence: A Modern Approach} (4th ed.). Pearson.

\bibitem{ref2}
Mitchell, T. M. (1997). \textit{Machine Learning}. McGraw-Hill Education.

\bibitem{ref3}
Tukey, J. W. (1977). \textit{Exploratory Data Analysis}. Addison-Wesley.

\bibitem{ref4}
Han, J., Pei, J., \& Tong, H. (2022). \textit{Data Mining: Concepts and Techniques} (4th ed.). Morgan Kaufmann/Elsevier.

\bibitem{ref5}
McKinney, W. (2017). \textit{Python for Data Analysis: Data Wrangling with Pandas, NumPy, and IPython} (2nd ed.). O'Reilly Media.

\bibitem{ref6}
Pedregosa, F., Varoquaux, G., Gramfort, A., Michel, V., Thirion, B., Grisel, O., Blondel, M., Prettenhofer, P., Weiss, R., Dubourg, V., Vanderplas, J., Passos, A., Cournapeau, D., Brucher, M., Perrot, M., \& Duchesnay, É. (2011). Scikit-learn: Machine Learning in Python. \textit{Journal of Machine Learning Research}, 12, 2825--2830.

\end{thebibliography}

\end{document}
